\documentclass{article}

\usepackage{amsmath,amssymb}
\usepackage{xurl}
\usepackage[hidelinks]{hyperref}
\setlength{\emergencystretch}{2em}

% Shared commands for notes in docs/paper-gaps/.

\DeclareUnicodeCharacter{2080}{\ifmmode{}_{0}\else\textsubscript{0}\fi}
\DeclareUnicodeCharacter{2081}{\ifmmode{}_{1}\else\textsubscript{1}\fi}
\DeclareUnicodeCharacter{2082}{\ifmmode{}_{2}\else\textsubscript{2}\fi}
\DeclareUnicodeCharacter{2099}{\ifmmode{}_{n}\else\textsubscript{n}\fi}

\newcommand{\Lean}{\textsc{Lean}}
\ExplSyntaxOn
\NewDocumentCommand{\leanid}{m}{
  \ifmmode
    \textsf{\detokenize{#1}}
  \else
    \group_begin:
    \urlstyle{sf}
    \tl_set:Nn \l_tmpa_tl {#1}
    \tl_replace_all:Nnn \l_tmpa_tl {\_} {_}
    \exp_args:NV \path \l_tmpa_tl
    \group_end:
  \fi
}
\ExplSyntaxOff
\providecommand{\tr}{\operatorname{tr}}
\newcommand{\tnleanIssueBase}{https://github.com/LionSR/TNLean/issues/}
\newcommand{\tnleanPrBase}{https://github.com/LionSR/TNLean/pull/}
\newcommand{\tnleanDocBase}{https://sirui-lu.com/TNLean/docs/}
\newcommand{\ghissue}[1]{\href{\tnleanIssueBase#1}{GitHub issue \##1}}
\newcommand{\ghpr}[1]{\href{\tnleanPrBase#1}{PR \##1}}
\newcommand{\doclink}[1]{\href{\tnleanDocBase#1.html}{\path{#1}}}

% Dirac notation and the identity operator, as in blueprint/src/macros/common.tex.
% \providecommand keeps note-local definitions authoritative.
\usepackage{dsfont}
\providecommand{\ket}[1]{|#1\rangle}
\providecommand{\bra}[1]{\langle #1|}
\providecommand{\braket}[2]{\langle #1 | #2 \rangle}
\providecommand{\ketbra}[2]{|#1\rangle\!\langle #2|}
\providecommand{\Id}{\mathds{1}}
\providecommand{\one}{\mathds{1}}
\providecommand{\C}{\mathbb{C}}
\providecommand{\MN}[1]{M_{#1}(\C)}
\providecommand{\MD}{M_D(\C)}

% External referencing for paper sources.  The build helper writes sanitized
% auxiliary files under build/paper-aux/ and excludes duplicated source labels.
\usepackage{xr}

% Import a generated paper auxiliary file with a caller-supplied label prefix.
% The candidate paths support builds from docs/paper-gaps/, the repository root,
% and build/paper-aux/ itself.
\newcommand{\importpaperaux}[2]{%
  \IfFileExists{../../build/paper-aux/#2.aux}{%
    \externaldocument[#1]{../../build/paper-aux/#2}%
  }{%
    \IfFileExists{build/paper-aux/#2.aux}{%
      \externaldocument[#1]{build/paper-aux/#2}%
    }{%
      \IfFileExists{#2.aux}{%
        \externaldocument[#1]{#2}%
      }{}%
    }%
  }%
}

% General external reference macros
\newcommand{\paperref}[2]{\ref{#1-#2}}
\newcommand{\papereqref}[2]{(\ref{#1-#2})}

% CPSV16 source (1606.00608 / MPDO-22-12-17-2.tex) external document and shortcuts
\importpaperaux{cpsv16-}{1606.00608/MPDO-22-12-17-2-tnlean}
\newcommand{\cpsvref}[1]{\paperref{cpsv16}{#1}}
\newcommand{\cpsveqref}[1]{\papereqref{cpsv16}{#1}}

% Verdict marker: \gapnote{<kind>}{<status>}. Kind is the policy
% classification (clarification, local-correction, scope-restriction,
% unfaithful, false-source, open-gap); status is open, wip (elimination
% underway), resolved, or historical. Machine-read by the site index;
% prints a verdict line.
\newcommand{\gapnote}[2]{\par\noindent\textbf{Verdict:} #1 (#2).\par}


\title{Schur Complement TFAE --- Contraction Gap}
\author{}
\date{2026-08-04}

\begin{document}
\maketitle

\section{Source statement}
Wolf, \textit{Quantum Channels \& Operations}, Theorem~5.2 (Block matrices and
Schur complements). Let $P$, $R$ be positive semi-definite complex matrices
and $Q$ a complex matrix of compatible dimensions. Then the following are
equivalent:
\begin{enumerate}
  \item the block matrix
    $\begin{pmatrix} P & Q \\ Q^\dagger & R \end{pmatrix}$
    is positive semi-definite;
  \item $\ker(R) \subseteq \ker(Q)$ and
    \begin{align}
      P \ge Q R^{+} Q^\dagger,
      \tag{Wolf, Theorem 5.2 (2)}
    \end{align}
    where $R^{+}$ is the pseudo-inverse (the inverse on the support);
  \item $\ker(R) \subseteq \ker(Q)$ and
    $\|P^{-1/2} Q R^{-1/2}\|_\infty \le 1$.
\end{enumerate}

\section{What is formalized}
The objects entering the statement are formalised: the block matrix itself,
its self-adjointness for Hermitian diagonal blocks, and the pseudoinverse
Schur complement $P - Q R^{+} Q^\dagger$.\footnote{Formal declarations:
\leanid{SchurComplement.blockMatrix}, \leanid{SchurComplement.blockMatrix_isHermitian},
and \leanid{SchurComplement.schurComplement} in
\doclink{TNLean/Channel/Schwarz/SchurComplement}.}
Conditions (1) and (2) are equivalent with the source's exact hypotheses:
the block matrix is PSD if and only if $\ker(R)\subseteq\ker(Q)$ and
$P-QR^+Q^\dagger$ is PSD.\footnote{Formal declaration:
\leanid{SchurComplement.blockMatrix_posSemidef_iff}.}
The forward implication combines the kernel inclusion
\leanid{SchurComplement.ker_subset_of_block_psd} with quadratic-form
minimisation at $y=-R^+Q^\dagger x$. The reverse implication uses the completed
square
\begin{align}
  \begin{pmatrix}x\\y\end{pmatrix}^{\!\dagger}
  \begin{pmatrix}P&Q\\Q^\dagger&R\end{pmatrix}
  \begin{pmatrix}x\\y\end{pmatrix}
  &=(R^+Q^\dagger x+y)^\dagger R(R^+Q^\dagger x+y)
    +x^\dagger(P-QR^+Q^\dagger)x.
\end{align}
This identity is formalized as
\leanid{SchurComplement.schur_complement_quadratic_form}. It rests on the
support absorption $Q\operatorname{supp}(R)=Q$ and the pseudoinverse identities
$RR^+=R^+R=\operatorname{supp}(R)$ and $(R^+)^\dagger=R^+$.

\section{What remains}
It remains to prove (2)$\Longleftrightarrow$(3). Algebraic manipulation with
$P^{1/2}$ and $R^{1/2}$ identifies $P \ge Q R^{+} Q^\dagger$ with the contraction
bound $\|P^{-1/2} Q R^{-1/2}\|_\infty \le 1$. The reverse route uses the
operator-norm characterization $\|L\| \le 1 \iff L L^\dagger \le I$ and the
support extensions of both inverse square roots. These operator-norm and
support-factorization statements are not yet formalized in the required form.

\section{Verdict}
Wolf Theorem~5.2 is partially formalized: conditions (1) and (2) are proved
equivalent without any additional hypothesis. Only the contraction equivalence
(2)$\Longleftrightarrow$(3) remains open, tracked in \ghissue{5501}.

\end{document}
